\section{Future Work}
\label{sec:future}

Of particular interest for future work is the scalability of Cornerstone regarding memory. As mentioned in \ref{subsec:local_dist} the factor by which memory decreases due to a doubling in GPU hardware is not close to $50\%$ as it would be in an ideal case, it was typically between $10\% - 20\%$ at the number of particles we tested. In order to efficiently reduce the memory consumption, which as we showed becomes the bottleneck, one would need to determine which structures within Cornerstone are causing scalability issues and optimize them where ever possible with an ideal, although unrealistic, target of reaching a $50\%$ reduction in particles when doubling the number of GPUs.

The method presented in \ref{subsec:vd_kf_opt} provides superior performance; however, it would need extensions to be more general and work for an arbitrary number of particle properties. Currently because the Radix sort is implemented via CUB the source code is unavailable, however, one could build a custom Radix Sort which sorts all properties within the kernel, achieving the same kernel fusion effect but with a pack of arrays rather than a single array using NVIDIA vector datatypes. This would be more generalizable than the vector datatype implementation because it can handle such arbitrary numbers of properties. NVIDIA also reportedly has plans to allow the CUB Radix sort to do byte-wise reordering which would also solve this problem. Additionally, a Thrust version could be implemented using zip iterators.

Octree's are structures built to service an algorithm whether it be collision detection, force calculations via Barnes Hut or Fast Multipole Method, graphics, or other applications. Understanding in what application contexts and under what conditions the tree build time is a bottleneck informs whether or not it is worth having a very optimized implementation of tree builds. Cornerstone in particular has a unique feature that affects performance being that the tree is rebalanced rather than having to rebuild this may be advantageous in some applications but cause unneeded complications in other problems.
